\section{Properties of determinants}

At this point we know how to compute determinants of small matrices, and we know a general method: Laplace expansion. But if we use Laplace expansion blindly, determinant calculations can become long and unpleasant. Mathematics is not only about having a method. It is also about knowing when a method can be simplified.

The determinant has several important properties. These properties are not just facts to memorize. They are tools for reading a matrix before calculating.

\subsection{Triangular matrices}

A square matrix is upper triangular if all entries below the main diagonal are zero. It is lower triangular if all entries above the main diagonal are zero.

For example,
\[
U=\begin{pmatrix}
2&-1&4\\
0&3&5\\
0&0&-2
\end{pmatrix}
\]
is upper triangular.

\begin{theorembox}
The determinant of a triangular matrix is the product of its diagonal entries.
\[
\det U=u_{11}u_{22}\cdots u_{nn}.
\]
\end{theorembox}

For the matrix above,
\[
\det U=2\cdot3\cdot(-2)=-12.
\]
This is much faster than expanding the determinant directly.

The reason is visible from Laplace expansion. If we expand a triangular matrix along a row or column with many zeros, the determinant keeps reducing to smaller triangular matrices. Eventually only diagonal entries remain.

\subsection{Swapping two rows}

If two rows of a matrix are exchanged, the determinant changes sign.

\begin{theorembox}
If \(B\) is obtained from \(A\) by swapping two rows, then
\[
\det B=-\det A.
\]
The same is true for swapping two columns.
\end{theorembox}

For a \(2\times2\) matrix this can be checked directly:
\[
\det\begin{pmatrix}a&b\\c&d\end{pmatrix}=ad-bc,
\]
while
\[
\det\begin{pmatrix}c&d\\a&b\end{pmatrix}=cb-da=-(ad-bc).
\]
The general case follows from the same alternating nature of the determinant.

This property explains why the sign pattern in cofactors matters. The determinant remembers orientation. Exchanging two rows reverses that orientation.

\subsection{Equal rows and proportional rows}

If two rows of a matrix are equal, then the determinant is zero.

Indeed, suppose two rows are equal. If we swap these two rows, the matrix does not change. Therefore the determinant should remain the same. But swapping two rows changes the sign of the determinant. Hence
\[
\det A=-\det A,
\]
so
\[
2\det A=0,
\]
and therefore
\[
\det A=0.
\]

This is a good example of a short argument. We did not compute the determinant. We used a property to force its value.

If two rows are proportional, the determinant is also zero. For example,
\[
A=\begin{pmatrix}
1&3&-2\\
2&6&-4\\
0&1&5
\end{pmatrix}
\]
has the second row equal to twice the first row. The determinant is zero because the first two rows do not contain independent information.

\begin{remarkbox}
A determinant equal to zero often means that some hidden dependence is present. Repeated rows and proportional rows are the easiest visible cases of this phenomenon.
\end{remarkbox}

\subsection{Multiplying a row by a number}

If one row of a matrix is multiplied by a number \(\lambda\), then the determinant is also multiplied by \(\lambda\).

\begin{theorembox}
If \(B\) is obtained from \(A\) by multiplying one row by \(\lambda\), then
\[
\det B=\lambda \det A.
\]
The same is true for multiplying one column by \(\lambda\).
\end{theorembox}

For example,
\[
\det\begin{pmatrix}
2&4\\
3&5
\end{pmatrix}=2\cdot5-4\cdot3=10-12=-2.
\]
If we multiply the first row by \(3\), we get
\[
\det\begin{pmatrix}
6&12\\
3&5
\end{pmatrix}=6\cdot5-12\cdot3=30-36=-6.
\]
The new determinant is three times the old one:
\[
-6=3(-2).
\]

This property is often used when simplifying determinants or when factoring common numbers out of rows or columns.

\subsection{Adding a multiple of one row to another}

One of the most useful properties is the following.

\begin{theorembox}
Adding a multiple of one row to another row does not change the determinant.
\end{theorembox}

This property is extremely important because it allows us to use row operations to simplify a determinant without changing its value. It is the determinant version of a powerful idea: transform the object into a simpler one while preserving the quantity we care about.

\begin{examplebox}
Compute
\[
\det\begin{pmatrix}
1&2&1\\
2&5&3\\
1&0&4
\end{pmatrix}.
\]
We simplify the matrix using row operations that do not change the determinant. Replace
\[
R_2\leftarrow R_2-2R_1,
\qquad
R_3\leftarrow R_3-R_1.
\]
Then
\[
\begin{pmatrix}
1&2&1\\
2&5&3\\
1&0&4
\end{pmatrix}
\longrightarrow
\begin{pmatrix}
1&2&1\\
0&1&1\\
0&-2&3
\end{pmatrix}.
\]
Now replace
\[
R_3\leftarrow R_3+2R_2.
\]
We get
\[
\begin{pmatrix}
1&2&1\\
0&1&1\\
0&0&5
\end{pmatrix}.
\]
This matrix is upper triangular, so its determinant is the product of the diagonal entries:
\[
\det A=1\cdot1\cdot5=5.
\]
\end{examplebox}

This calculation is much more efficient than expanding immediately. It also demonstrates how determinant properties support intelligent computation.

\subsection{The determinant of a product}

Another fundamental property is
\[
\det(AB)=\det(A)\det(B),
\]
for square matrices \(A\) and \(B\) of the same size.

This formula says that determinants multiply under matrix multiplication. If one matrix scales area or volume by one factor and another matrix scales it by another factor, then doing both transformations multiplies the scale factors.

One important consequence is immediate. If \(A\) has determinant zero, then for every compatible square matrix \(B\),
\[
\det(AB)=0.
\]
A matrix that has already collapsed structure cannot be repaired by multiplying by another matrix on one side.

\subsection{How to use the properties}

The properties of determinants are not separate tricks. They form a strategy.

\begin{summarybox}
Before computing a determinant, ask:
\begin{itemize}
\item Is the matrix triangular?
\item Are two rows or columns equal or proportional?
\item Can I create zeros using row operations that preserve the determinant?
\item If I swap rows, did I change the sign?
\item If I multiply a row by a number, did I multiply the determinant by that number?
\item Is Laplace expansion now easier after simplification?
\end{itemize}
\end{summarybox}

A determinant should not be expanded mechanically unless there is no better path. Often the best computation begins with looking at the matrix and asking what structure is already visible.
